% Sources verifiees pour les valeurs numeriques non triviales de ce fichier :
%
% NOTE DE COMPLEMENTARITE (2026-09-02) : ce fichier a ete revu pour ne plus dupliquer les
% Questions Types de la fiche fiches/T12_Machine_Synchrone.tex. Les QT de la fiche sont :
%   QT1 = 2024 E2 appl. Q44-46 (f0 = 185,6 Hz + couplage triangle, Urbanloop)
%   QT2 = 2025 E2 appl. Q18 (definition du servomoteur, KUKA)
%   QT3 = 2023 E2 disc. (rendement de la chaine de puissance, Bakus)
% Aucun exercice de ce fichier ne reprend ces trois questions.
%
% - Ex1 (docTech, choix du servomoteur brushless d'axe 6 KUKA) : CAPET 2025, Epreuve ecrite
%   disciplinaire appliquee, questions 11, 12 et 13.
%   * VIENT DU SUJET (textuellement present dans la conversion Markdown) :
%     - "Pour les questions suivantes, Jeq = 0,0118 kg.m2 et le couple resistif Cr est neglige"
%       (E2_Appliquee_Commune_2025.md, l.310) ;
%     - le contenu du DT2 "Extrait de la gamme des robots KUKA" (l.776-784, texte de l'image
%       extrait par Marker) : KR 300 2700 -> masse maximale 370 kg, Cm max = 8,21 N.m ;
%       KR 210 R3100 -> masse maximale 210 kg, Cm max = 3,83 N.m ;
%     - l'acceleration angulaire a = 431,1 rad/s2 et la masse de 181 kg de l'ensemble
%       "structure + ecrans" (Rapport_Jury_2025.md, l.~ Q10, Q12 et Q16 : "a=517,31/1,2 donc
%       a = 431,1 rad.s-2", "Le couple necessaire pour mouvoir la masse de 181 kg est de
%       5,086 N.m", "Cm = 5,086 N.m") ;
%     - Q13 corrige officiel : mmax = 329,9 kg (Rapport_Jury_2025.md, Q13).
%   * RECONSTITUE : la mise en forme tabulaire de l'extrait DT2 (le DT original est une image
%     "picture [448 x 619] intentionally omitted" ; seules les valeurs listees ci-dessus sont
%     reprises, aucune valeur ajoutee). La ligne "figure 12" donnant a = 431,1 rad/s2 est
%     fournie directement en donnee d'enonce, la courbe n'etant pas convertie.
% - Ex2 (quadrants de fonctionnement de la MS sur le profil de vitesse) : CAPET 2024, Epreuve
%   ecrite disciplinaire appliquee, question 44 -- V0 = 60 km/h, deceleration aC = -1,4 m/s2
%   (E2_Appliquee_Commune_2024.md, l.470-472). Le corrige officiel de Q44 est une image omise
%   par la conversion (Rapport_Jury_2024.md, "Q44." l.~880) : le raisonnement quadrant par
%   quadrant est ici reconstruit a partir du seul profil de vitesse, sans citation attribuee.
% - Ex3 (resolveur + chaine d'acquisition de position) : CAPET 2025, Epreuve ecrite
%   disciplinaire appliquee, questions 25, 26 et 27 -- toutes les donnees (resolveur, plage
%   [-5 V ; +5 V], debattement 288 deg code sur 24 bits, critere q = dv/10) sont dans le sujet
%   (E2_Appliquee_Commune_2025.md, l.374-394) et les resultats (dtheta = 1,72e-5 deg,
%   q ~ 6e-8 V, N ~ 1,6e8, 28 bits) dans le corrige officiel (Rapport_Jury_2025.md, Q23, Q27).
% - Ex4 (reversibilite pont PD2 + dimensionnement supercondensateurs + VERIFICATION du
%   dimensionnement) : CAPET 2024, questions 48, 49, 51, 52 et 53 -- Nsc=5, Ceq=13 F,
%   UCeqmax=81 V, formulation "unidirectionnel en courant et en tension", energie recuperee
%   -139 kJ a 80 % et conclusion "L'agencement propose a la question 51 ne convient pas car
%   UCeqf < UCeqmax" avec ses trois solutions : tous du corrige officiel
%   (Rapport_Jury_2024.md, Q48-Q53) ; 65 F / 16,2 V et U0 = 72 V du sujet
%   (E2_Appliquee_Commune_2024.md, l.602-608 et l.510).
% - Ex5 (diagramme de Kapp) : exercice inedit, theorie generale de la machine synchrone (non
%   controversee), aucune verification web necessaire.
%
% ANCIEN Ex3 (chaine de puissance Bakus 2023) SUPPRIME car doublon de la QT3 de la fiche.
% Rappel pour toute reecriture future : le sujet 2023 ne donne QUE le rendement global de
% 63 % ; il n'existe AUCUNE decomposition officielle en etaCS/etaMS/etaRed, et le facteur
% 98 % du sujet concerne l'onduleur du hangar, pas la chaine de traction. Ne jamais
% reintroduire un produit "0,98 x 0,80 x 0,80" presente comme venant du sujet.
\expandafter\ifx\csname CAPETINCLUDEMODE\endcsname\relax
\documentclass[11pt,a4paper]{article}
% preamble.tex — packages et commandes communes a toutes les fiches CAPET SII IE
\input{glyphtounicode}
\pdfgentounicode=1
\usepackage{cmap}
\usepackage[T1]{fontenc}
\usepackage[utf8]{inputenc}
\usepackage[french]{babel}
% babel[french] active a \begin{document} des raccourcis de ponctuation (:;!?)
% qui cassent la syntaxe to[...] de circuitikz (erreur "+ or - expected") ;
% on les desactive, ce qui ne change que l'espacement fin avant ces signes.
\AtBeginDocument{\shorthandoff{:;!?}}
\usepackage[a4paper, top=1.6cm, bottom=1.6cm, left=1.4cm, right=1.4cm]{geometry}
\setlength{\headheight}{14.5pt}
\usepackage{amsmath, amssymb, amsfonts}
\usepackage{siunitx}
\sisetup{locale=FR, output-decimal-marker={,}, inter-unit-product=\ensuremath{{}\cdot{}},
         group-minimum-digits=4}
% Unites non fournies par siunitx, declarees ici une fois pour toutes (fiches + exercices).
% Les gardes \ifdefined evitent tout conflit avec une redeclaration locale restee dans un
% fichier : ne pas retirer les gardes, ne pas redeclarer ces unites dans les fichiers.
\DeclareSIUnit{\jour}{j}
\ifdefined\var\else\DeclareSIUnit{\var}{var}\fi   % puissance reactive
\ifdefined\vad\else\DeclareSIUnit{\vad}{VAD}\fi   % puissance deformante
\ifdefined\tr\else\DeclareSIUnit{\tr}{tr}\fi      % tour (tr/min)
\usepackage{array, booktabs, multirow, tabularx}
\usepackage{xcolor}
\usepackage[most, breakable]{tcolorbox}
\usepackage{multicol}
\usepackage{enumitem}
\usepackage{mathtools}
\usepackage{lmodern}
% Refonte visuelle 2026 (direction "moderne coloré") : Lato pour le texte,
% maths sans-serif assorties via newtxsf (charge apres amsmath/amssymb).
\usepackage[default]{lato}
\usepackage{newtxsf}
\usepackage{microtype}
\usepackage{fancyhdr}
\usepackage{tikz}
\usetikzlibrary{positioning,arrows.meta}
\usepackage{circuitikz}
\usepackage{lastpage}

% ── Couleurs par domaine (palette modernisée 2026) ────────────────────
\definecolor{colElec}{RGB}{29, 78, 216}        % Électrotechnique — bleu vif
\definecolor{colElecLight}{RGB}{219, 234, 254} % fond bleu clair
\definecolor{colAuto}{RGB}{21, 128, 61}        % Automatique — vert
\definecolor{colMeca}{RGB}{194, 65, 12}        % Mécanique — orange
\definecolor{colTherm}{RGB}{185, 28, 28}       % Thermique — rouge
\definecolor{colQT}{RGB}{232, 244, 233}        % (hérité, plus utilisé par questiontype)
\definecolor{colJury}{RGB}{255, 248, 225}      % fond jury (crème)
\definecolor{colJuryFrame}{RGB}{202, 138, 4}   % bordure/bandeau jury (ambre)
\definecolor{colGrayLight}{RGB}{245, 245, 247} % fond gris neutre
\definecolor{colSeparator}{RGB}{180, 180, 180} % gris séparateur
\definecolor{colExo}{RGB}{8, 145, 178}          % Banque d'exercices — cyan (distinct des fiches)
% ── Encre et gris de la refonte ───────────────────────────────────────
\definecolor{colInk}{RGB}{23, 32, 42}          % texte courant (quasi-noir)
\definecolor{colSlate}{RGB}{83, 98, 118}       % gris ardoise (méta, en-têtes)
\definecolor{colRule}{RGB}{203, 213, 225}      % filets clairs
\definecolor{colJuryBg}{RGB}{255, 248, 225}    % fond crème jury
\definecolor{colPresence}{RGB}{83, 98, 118}    % bloc 5 (présence dans les sujets)
\AtBeginDocument{\color{colInk}}

% ── Style encadré exercice + corrigé (banque d'exercices) ─────────────
\tcbset{
  exoBox/.style={
    breakable,
    enhanced,
    colback=colExo!8!white,
    colframe=colExo,
    colbacktitle=colExo,
    coltitle=white,
    fonttitle=\bfseries\small,
    left=4pt, right=4pt, top=3pt, bottom=3pt,
    before skip=5pt, after skip=5pt,
  }
}

% ── Bascule enonce / corrige (banque d'exercices) ─────────────────────
% \CAPETCORRIGE est defini en ligne de commande (pdflatex "\def\CAPETCORRIGE{1}\input{...}")
% pour compiler la variante corriges d'un meme fichier source ; par defaut (non defini),
% la compilation produit la variante enonces. \CAPETSHOW{enonce}{corrige} choisit la branche.
%
% TROIS MODES (refonte 2026-09-20, manuel de revision) :
%   - aucun drapeau   -> imprime l'enonce            (variante enonces, INCHANGE)
%   - \CAPETCORRIGE   -> imprime le corrige          (variante corriges, INCHANGE)
%   - \CAPETMANUEL    -> imprime l'enonce ET met le corrige de cote pour l'annexe
%
% Le mode manuel stocke chaque corrige dans une macro (\csgdef), sans ecriture
% fichier ni re-tokenisation : le contenu (circuitikz, accents, %) traverse intact.
% L'annexe les rejoue ensuite dans l'ordre via \CAPETANNEXE.
\usepackage{etoolbox}

\newcounter{CAPETexo}          % numero continu d'exercice dans le manuel
\setcounter{CAPETexo}{0}
\newcounter{CAPETannexe}       % compteur de relecture en annexe

% Titre courant de la notion en cours (pose par \CAPETNOTION), sert d'intertitre annexe
\newcommand{\CAPETnotioncode}{}
\newcommand{\CAPETnotiontitre}{}

% \CAPETNOTION{code}{titre} : ouvre une notion dans le manuel (memorise le libelle,
% et marque le rang du premier exercice de la notion pour l'intertitre d'annexe).
\newcommand{\CAPETNOTION}[2]{%
  \renewcommand{\CAPETnotioncode}{#1}%
  \renewcommand{\CAPETnotiontitre}{#2}%
  \csgdef{CAPETnotionstart\the\numexpr\value{CAPETexo}+1\relax}{#1 -- #2}%
}

% ── \degree robuste ───────────────────────────────────────────────────
% Certains fichiers d'exercices utilisent \degree. Ni le preambule ni siunitx
% ne le definissent, et \providecommand local ne suffit pas : sous hyperref,
% \degree existe deja comme commande factice (\ERROR), donc \providecommand
% s'abstient et l'appel echoue a la compilation du manuel (erreur invisible
% en compilation isolee). On le (re)definit donc inconditionnellement ici,
% apres le chargement des packages, pour tout le corpus.
\AtBeginDocument{\providecommand{\degree}{}\renewcommand{\degree}{\ensuremath{^\circ}}}

% Libelle decoratif : les fichiers d'exercices utilisent aussi \CAPETSHOW dans
% leur \fancyhead et leur bloc de titre, pour afficher "(enonces)" ou "(corriges)".
% Ces appels-la ne sont PAS des exercices : en mode manuel ils ne doivent ni
% consommer un numero, ni alimenter l'annexe. On les reconnait a leur premier
% argument, strictement egal au libelle ci-dessous.
\newcommand{\CAPETLIBELLEENONCES}{(\'enonc\'es)}

\ifdefined\CAPETMANUEL
  % --- Mode manuel : enonce en place, corrige differe vers l'annexe ---
  \newcommand{\CAPETSHOW}[2]{%
    \def\CAPETargun{#1}%
    \ifx\CAPETargun\CAPETLIBELLEENONCES
      #1%            appel decoratif : rendu tel quel, hors numerotation
    \else
      \CAPETSHOWEXO{#1}{#2}%
    \fi
  }
  \newcommand{\CAPETSHOWEXO}[2]{%
    \stepcounter{CAPETexo}%
    \csgdef{CAPETcorrige\theCAPETexo}{#2}%
    \edef\CAPETtmpnum{\theCAPETexo}%
    #1%
    \par\vspace{2pt}%
    {\footnotesize\itshape\color{colExo}%
     \hyperref[capet:cor:\CAPETtmpnum]{Corrige de l'exercice \CAPETtmpnum\ $\rightarrow$ p.~\pageref{capet:cor:\CAPETtmpnum}}%
     \label{capet:eno:\CAPETtmpnum}}%
  }
  % \CAPETANNEXE : rejoue tous les corriges captures, dans l'ordre, avec intertitres
  \newcommand{\CAPETANNEXE}{%
    \setcounter{CAPETannexe}{0}%
    \loop\ifnum\value{CAPETannexe}<\value{CAPETexo}%
      \stepcounter{CAPETannexe}%
      \edef\CAPETcurnum{\theCAPETannexe}%
      \ifcsname CAPETnotionstart\CAPETcurnum\endcsname
        \par\medskip
        {\normalsize\bfseries\color{colExo}\csuse{CAPETnotionstart\CAPETcurnum}}%
        \par\nobreak\smallskip
      \fi
      \begin{tcolorbox}[exoBox, breakable,
        title={Corrige de l'exercice \CAPETcurnum\
               \hfill\normalfont\footnotesize
               \hyperref[capet:eno:\CAPETcurnum]{enonce p.~\pageref{capet:eno:\CAPETcurnum}}}]%
      \label{capet:cor:\CAPETcurnum}%
      \csuse{CAPETcorrige\CAPETcurnum}%
      \end{tcolorbox}%
    \repeat
  }
\else
  % --- Modes historiques, comportement strictement inchange ---
  \newcommand{\CAPETSHOW}[2]{\ifdefined\CAPETCORRIGE #2\else #1\fi}
  \newcommand{\CAPETANNEXE}{}
\fi

% ── Difficulté 1 à 3 (étoiles pleines, amssymb) ───────────────────────
\newcommand{\difficulte}[1]{%
  \ifcase#1\relax
  \or $\bigstar$\;niveau 1/3%
  \or $\bigstar\bigstar$\;niveau 2/3%
  \or $\bigstar\bigstar\bigstar$\;niveau 3/3%
  \fi
}

% ── Style encadré extrait de document technique (banque d'exercices) ──
\tcbset{
  docTech/.style={
    breakable,
    enhanced,
    colback=colGrayLight,
    colframe=colSeparator,
    fonttitle=\bfseries\small\itshape,
    coltitle=black,
    colbacktitle=colGrayLight,
    left=4pt, right=4pt, top=3pt, bottom=3pt,
    before skip=4pt, after skip=4pt,
  }
}

% ── Style tcolorbox de section principale ────────────────────────────
% Refonte 2026 : bandeau titre coloré arrondi + fond très pâle, sans cadre.
\tcbset{
  ficheSection/.style={
    breakable, enhanced jigsaw,
    frame hidden, boxrule=0pt, arc=1.8mm,
    colback=#1!4!white,
    colbacktitle=#1,
    coltitle=white,
    fonttitle=\sffamily\bfseries,
    left=7pt, right=7pt, top=5pt, bottom=5pt,
    toptitle=3pt, bottomtitle=3pt,
    before skip=10pt, after skip=8pt,
  }
}

% ── Style encadré question type ──────────────────────────────────────
% Carte blanche, bandeau sombre, liseré très clair.
\tcbset{
  questiontype/.style={
    breakable, enhanced jigsaw,
    colback=white,
    frame hidden, boxrule=0pt, arc=1.8mm,
    colbacktitle=colElec!78!black,
    coltitle=white,
    fonttitle=\sffamily\bfseries\small,
    left=7pt, right=7pt, top=4pt, bottom=4pt,
    toptitle=2.5pt, bottomtitle=2.5pt,
    before skip=8pt, after skip=8pt,
    borderline={0.6pt}{0pt}{colElec!12!white},
  }
}

% ── Style encadré jury ───────────────────────────────────────────────
% Bandeau ambre, fond crème, sans cadre.
\tcbset{
  jury/.style={
    breakable, enhanced jigsaw,
    colback=colJuryBg,
    frame hidden, boxrule=0pt, arc=1.8mm,
    colbacktitle=colJuryFrame,
    coltitle=white,
    fonttitle=\sffamily\bfseries,
    left=7pt, right=7pt, top=4pt, bottom=4pt,
    toptitle=3pt, bottomtitle=3pt,
    before skip=10pt, after skip=8pt,
  }
}

% ── En-tête / pied de page communs (refonte 2026) ─────────────────────
% #1 code fiche, #2 titre court, #3 couleur domaine
\newcommand{\ficheHeaders}[3]{%
  \fancyhead[L]{\small\sffamily{\bfseries\textcolor{#3}{#1}}\hspace{0.7em}\textcolor{colSlate}{#2}}%
  \fancyhead[R]{\small\sffamily\textcolor{colSlate}{CAPET SII IE\, --- \,Session 2026}}%
  \fancyfoot[C]{\small\sffamily\textcolor{colSlate}{\thepage\,/\,\pageref{LastPage}}}%
  \fancyhead[C]{}\fancyfoot[L]{}\fancyfoot[R]{}%
}

% ── Bloc en-tête de fiche (BLOC 0, refonte 2026) ──────────────────────
% #1 code  #2 titre  #3 couleur domaine  #4 domaine  #5 ligne présence  #6 sous-thèmes
\newcommand{\ficheEntete}[6]{%
  \begin{tcolorbox}[enhanced, breakable,
    frame hidden, boxrule=0pt, arc=2.5mm,
    interior style={left color=#3, right color=#3!75!black},
    left=9pt, right=9pt, top=7pt, bottom=7pt,
    before skip=0pt, after skip=6pt]
    {\color{white}\sffamily\bfseries\LARGE #1\quad #2}\\[4pt]
    {\color{#3!22!white}\sffamily\small \MakeUppercase{#4}\quad\textbar\quad #5}\\[1pt]
    {\color{#3!35!white}\sffamily\footnotesize Sous-thèmes : #6}
  \end{tcolorbox}%
}

% ── Mise en page ─────────────────────────────────────────────────────
\pagestyle{fancy}
\fancyhf{}
\renewcommand{\headrulewidth}{0.5pt}
\renewcommand{\footrulewidth}{0pt}
\AtBeginDocument{%
  \renewcommand{\headrule}{{\color{colRule}\hrule width\headwidth height 0.5pt}}%
}

\setlist[itemize]{leftmargin=1.2em, itemsep=1pt, topsep=2pt, parsep=0pt}
\setlist[enumerate]{leftmargin=1.2em, itemsep=3pt, topsep=2pt, parsep=0pt}

\setlength{\parskip}{2pt}
\setlength{\columnsep}{10pt}

\def\CAPETBEGINFILE{\begin{document}}
\def\CAPETENDFILE{\end{document}}
\else
\def\CAPETBEGINFILE{}
\def\CAPETENDFILE{}
\fi

% Macro locale : numero de phase entoure (profil de vitesse de l'exercice 2).
\providecommand{\circled}[1]{\tikz[baseline=(c.base)]{\node[draw, circle, inner sep=1pt, minimum size=1.3em, line width=0.4pt] (c) {\footnotesize #1};}}

\fancyhead[L]{\textbf{\textcolor{colExo}{T12}} - Exercices - Machine synchrone et brushless (PMSM / BLDC)}
\fancyhead[R]{\small\textit{Banque d'exercices \CAPETSHOW{(\'enonc\'es)}{(corrig\'es)}}}
\fancyfoot[C]{\small\thepage\,/\,\pageref{LastPage}}

\CAPETBEGINFILE

\begin{tcolorbox}[
  enhanced, breakable,
  colback=colExo,
  colframe=colExo,
  left=6pt, right=6pt, top=5pt, bottom=5pt,
  before skip=0pt, after skip=4pt
]
  {\color{white}\bfseries\large T12 \quad Machine synchrone et brushless (PMSM / BLDC)}\\[2pt]
  {\color{white}\small Banque d'exercices -- 5 exercices \CAPETSHOW{(\'enonc\'es)}{(corrig\'es)}}
\end{tcolorbox}

\begin{tcolorbox}[exoBox,
  title={Exercice 1 -- \difficulte{1} \hfill\normalfont\footnotesize 2025, \'Epreuve \'ecrite disciplinaire appliqu\'ee (Sc\`ene robotis\'ee KUKA), questions 11, 12 et 13 -- \textbf{lecture de document technique}}
]
\CAPETSHOW{
Une soci\'et\'e de production souhaite suspendre une structure porteuse d'\'ecrans g\'eants \`a l'axe~6 d'un robot industriel KUKA. Cet axe est entra\^in\'e par un \textbf{servomoteur brushless}. Il faut choisir, dans la gamme du constructeur, le robot dont le moteur d'axe~6 fournit le couple n\'ecessaire.

Le moment d'inertie \'equivalent ramen\'e \`a l'arbre du moteur d'axe~6 vaut $J_{eq} = \SI{0,0118}{\kilogram\meter\squared}$ et le couple r\'esistif $C_r$ est n\'eglig\'e. Le relev\'e de la vitesse du moteur de l'axe~6 donne une acc\'el\'eration angulaire $a = \SI{431,1}{\radian\per\second\squared}$. L'ensemble \guillemotleft\ structure $+$ \'ecrans \guillemotright{} \`a embarquer a une masse de $\SI{181}{\kilogram}$.

\begin{tcolorbox}[docTech, title={Document technique DT2 -- Extrait de la gamme des robots KUKA}]
\centering\small
\begin{tabular}{@{}lcc@{}}
\toprule
\textbf{D\'esignation} & \textbf{Masse maximale} & \textbf{Couple max.\ moteur d'axe 6} \\
 & (charge admissible) & $C_{m\,max}$ \\
\midrule
KR 300 R2700  & \SI{370}{\kilogram} & \SI{8,21}{\newton\meter} \\
KR 210 R3100  & \SI{210}{\kilogram} & \SI{3,83}{\newton\meter} \\
\bottomrule
\end{tabular}

\smallskip
\raggedright\footnotesize\itshape Le diagramme de charge associ\'e \`a chaque r\'ef\'erence donne, pour un robot donn\'e, la masse maximale admissible en fonction de la position du centre de gravit\'e de la charge ($L_{xy}$, $L_z$) par rapport \`a la bride outil.
\end{tcolorbox}

\begin{enumerate}[leftmargin=1.5em]
  \item \textbf{Relever dans le document technique DT2} les couples maximaux disponibles sur le moteur d'axe~6 de chacune des deux r\'ef\'erences propos\'ees, ainsi que les masses maximales admissibles associ\'ees.
  \item Le principe fondamental de la dynamique appliqu\'e \`a un syst\`eme en rotation autour d'un axe fixe s'\'ecrit $\sum C = J_{eq}\cdot a$. Calculer le couple moteur $C_m$ n\'ecessaire pour mouvoir la structure embarqu\'ee.
  \item \`A partir du DT2, choisir en le justifiant le robot r\'epondant \`a ce contexte d'utilisation.
  \item Pour le robot retenu, le moment d'inertie \'equivalent s'exprime $J_{eq} = J_{mot6} + 1{,}32\,m_{max}\,r^2$ avec $J_{mot6} = \SI{0,003}{\kilogram\meter\squared}$ et $r = 0{,}006\,07$ (rapport de r\'eduction). D\'eterminer la masse maximale $m_{max}$ que le robot peut r\'eellement embarquer en tenant compte des effets dynamiques, et conclure.
\end{enumerate}

}{
\textbf{Corrig\'e.}

\begin{enumerate}[leftmargin=1.5em]
  \item Extraction directe du DT2 :
  \begin{itemize}[leftmargin=1.5em]
    \item KR 300 R2700 : $C_{m\,max} = \SI{8,21}{\newton\meter}$, masse maximale $\SI{370}{\kilogram}$ ;
    \item KR 210 R3100 : $C_{m\,max} = \SI{3,83}{\newton\meter}$, masse maximale $\SI{210}{\kilogram}$.
  \end{itemize}
  Les deux r\'ef\'erences admettent une masse sup\'erieure aux $\SI{181}{\kilogram}$ de la charge : le crit\`ere statique de masse ne suffit donc pas \`a d\'epartager, c'est le \emph{couple} d'axe~6 qui va trancher.

  \item Le couple r\'esistif \'etant n\'eglig\'e, le PFD en rotation se r\'eduit \`a $C_m = J_{eq}\cdot a$ :
  \[ C_m = 0{,}011\,8 \times 431{,}1 \approx \SI{5,086}{\newton\meter} \]

  \item Comparaison au DT2 :
  \begin{itemize}[leftmargin=1.5em]
    \item KR 210 R3100 : $C_{m\,max} = \SI{3,83}{\newton\meter} < \SI{5,086}{\newton\meter}$ $\rightarrow$ \textbf{sous-dimensionn\'e}, le servomoteur ne pourrait pas fournir l'acc\'el\'eration demand\'ee ;
    \item KR 300 R2700 : $C_{m\,max} = \SI{8,21}{\newton\meter} > \SI{5,086}{\newton\meter}$ $\rightarrow$ \textbf{convient}.
  \end{itemize}
  Le choix se porte donc sur le \textbf{KUKA KR 300 R2700}.

  \item On \'ecrit la condition limite $C_{m\,max} = J_{eq}\cdot a$, soit $J_{eq} = C_{m\,max}/a$ :
  \[ 0{,}003 + 1{,}32\,m_{max}\,(0{,}006\,07)^2 = \dfrac{8{,}21}{431{,}1} = \SI{0,019\,04}{\kilogram\meter\squared} \]
  \[ 1{,}32\,m_{max}\times 3{,}684\times10^{-5} = 0{,}019\,04 - 0{,}003 = 0{,}016\,04 \]
  \[ m_{max} = \dfrac{0{,}016\,04}{4{,}863\times10^{-5}} \approx \SI{330}{\kilogram} \]
  (le corrig\'e officiel donne $m_{max} = \SI{329,9}{\kilogram}$.)
\end{enumerate}

\textbf{Conclusion.} La masse maximale r\'eellement embarquable sur l'axe~6 du KR 300 R2700, $\SI{330}{\kilogram}$, est fix\'ee par le \emph{couple} du servomoteur brushless et non par la charge nominale $\SI{370}{\kilogram}$ affich\'ee au catalogue : c'est bien le crit\`ere dynamique qui dimensionne. Avec $\SI{181}{\kilogram}$ embarqu\'es, la marge est de $(330-181)/181 \approx \SI{82}{\percent}$, ce qui valide largement l'exigence. Attention toutefois : ce calcul ne valide que l'axe~6, suppos\'e le plus p\'enalisant -- le diagramme de charge du DT2 doit \^etre utilis\'e pour v\'erifier tous les axes en fonction de la position du centre de gravit\'e de la charge.
}
\end{tcolorbox}

\begin{tcolorbox}[exoBox,
  title={Exercice 2 -- \difficulte{1} \hfill\normalfont\footnotesize 2024, \'Epreuve \'ecrite disciplinaire appliqu\'ee (Urbanloop), question 44}
]
\CAPETSHOW{
Chaque roue de la capsule Urbanloop est motoris\'ee par une machine synchrone \`a aimants permanents entra\^inant directement la roue. Sur une portion horizontale et rectiligne du parcours, le profil de vitesse d'une capsule entre deux stations est le suivant : d\'emarrage depuis l'arr\^et, phase d'acc\'el\'eration jusqu'\`a la vitesse de croisi\`ere $V_0 = \SI{60}{\kilo\meter\per\hour}$, palier \`a vitesse constante, puis phase de d\'ec\'el\'eration ($a_C = \SI{-1,4}{\meter\per\second\squared}$) jusqu'\`a l'arr\^et en station. La capsule circule toujours dans le m\^eme sens.

\begin{center}
\begin{tikzpicture}[scale=1.0, >=Stealth]
  \draw[-latex] (0,0) -- (8.5,0) node[right]{$t$};
  \draw[-latex] (0,0) -- (0,2.4) node[above]{$v$};
  \draw[dashed, gray] (0,1.8) -- (6.2,1.8); \node[left] at (0,1.8) {\small $V_0$};
  \draw[very thick, colExo] (0,0) -- (2.0,1.8) -- (5.0,1.8) -- (6.2,0);
  \node[below] at (1.0,0) {\small \circled{1}};
  \node[below] at (3.5,0) {\small \circled{2}};
  \node[below] at (5.6,0) {\small \circled{3}};
\end{tikzpicture}

{\footnotesize Profil de vitesse entre deux stations : \circled{1} acc\'el\'eration, \circled{2} vitesse constante, \circled{3} freinage.}
\end{center}

Le plan couple-vitesse $(\Omega, C)$ de la machine est d\'ecoup\'e en quatre quadrants, la vitesse $\Omega$ en abscisse et le couple $C$ en ordonn\'ee.

\begin{enumerate}[leftmargin=1.5em]
  \item Pour chacune des trois phases du trajet, pr\'eciser le signe de la vitesse de rotation $\Omega$ de la machine et le signe du couple $C$ qu'elle d\'eveloppe.
  \item En d\'eduire, pour chaque phase, le quadrant de fonctionnement occup\'e et le mode de fonctionnement de la machine synchrone (moteur ou g\'en\'erateur).
  \item En d\'eduire l'exigence que ce profil impose au convertisseur qui alimente la machine synchrone.
\end{enumerate}

}{
\textbf{Corrig\'e.}

\begin{enumerate}[leftmargin=1.5em]
  \item La capsule circule toujours dans le m\^eme sens : la vitesse de rotation des moteurs-roues est donc \textbf{toujours positive} sur tout le trajet ($\Omega > 0$, et $\Omega = 0$ aux deux extr\'emit\'es).
  \begin{itemize}[leftmargin=1.5em]
    \item \circled{1} \emph{Acc\'el\'eration} : la machine doit fournir le couple qui acc\'el\`ere la capsule et vaincre les frottements $\rightarrow$ $C > 0$.
    \item \circled{2} \emph{Vitesse constante} : le couple compense exactement la r\'esistance a\'erodynamique $\rightarrow$ $C > 0$ (plus faible qu'en phase \circled{1}).
    \item \circled{3} \emph{Freinage} : la machine doit s'opposer au mouvement pour imposer $a_C = \SI{-1,4}{\meter\per\second\squared}$ $\rightarrow$ le couple change de signe, $C < 0$.
  \end{itemize}

  \item Avec $\Omega > 0$ dans tous les cas, seuls deux quadrants sont parcourus :
  \begin{center}
  \begin{tabular}{@{}llll@{}}
  \toprule
  \textbf{Phase} & $\Omega$ & $C$ & \textbf{Quadrant / mode} \\
  \midrule
  \circled{1} Acc\'el\'eration      & $>0$ & $>0$ & Quadrant 1 -- \textbf{moteur} \\
  \circled{2} Vitesse constante     & $>0$ & $>0$ & Quadrant 1 -- \textbf{moteur} \\
  \circled{3} Freinage              & $>0$ & $<0$ & Quadrant 4 -- \textbf{g\'en\'erateur} \\
  \bottomrule
  \end{tabular}
  \end{center}
  Le produit $P = C\,\Omega$ est positif en \circled{1} et \circled{2} (la machine \emph{re\c coit} de la puissance \'electrique : fonctionnement moteur) et n\'egatif en \circled{3} (la machine \emph{fournit} de la puissance \'electrique : fonctionnement g\'en\'eratrice, c'est le freinage r\'ecup\'eratif).

  \item Le fonctionnement occupant les quadrants 1 et 4, la machine change de signe de couple sans changer de sens de rotation : c\^ot\'e continu, le courant doit pouvoir \textbf{s'inverser} alors que la tension du bus reste positive. Le convertisseur doit donc \^etre \textbf{r\'eversible en courant} (deux quadrants), ce que r\'ealise nativement l'onduleur de tension \`a transistors et diodes antiparall\`eles.
\end{enumerate}

\textbf{Conclusion.} Le trajet entre deux stations n'exige que deux quadrants sur quatre, mais le quadrant~4 (freinage) impose imp\'erativement une r\'eversibilit\'e en courant de la cha\^ine d'alimentation. C'est pr\'ecis\'ement ce point qui va poser probl\`eme en amont de l'onduleur (voir exercice~4) : si l'\'etage qui alimente le bus continu n'est pas lui-m\^eme r\'eversible, l'\'energie de freinage n'a nulle part o\`u aller.
}
\end{tcolorbox}

\begin{tcolorbox}[exoBox,
  title={Exercice 3 -- \difficulte{2} \hfill\normalfont\footnotesize 2025, \'Epreuve \'ecrite disciplinaire appliqu\'ee (Sc\`ene robotis\'ee KUKA), questions 25, 26 et 27}
]
\CAPETSHOW{
Le servomoteur brushless de l'axe~3 d'un robot KUKA KR QUANTEC est asservi en position gr\^ace \`a un capteur de position angulaire de type \textbf{r\'esolveur}, dont les deux signaux de sortie correspondent au cosinus et au sinus de la position angulaire de l'axe du moteur.

L'axe en sortie du servomoteur dispose d'un d\'ebattement angulaire de $288^\circ$ et re\c coit des commandes de position cod\'ees sur $24$ bits, le d\'ebattement \'etant enti\`erement parcouru entre les deux positions extr\^emes.

Un circuit de traitement convertit l'information de position angulaire en une tension image valant $\SI{-5}{\volt}$ lorsque l'axe est en $\theta_{min}$ et $\SI{+5}{\volt}$ lorsqu'il est en $\theta_{max}$. On note $\delta v$ la variation de cette tension image lorsque le moteur tourne d'un angle $\delta\theta$ (le pas angulaire commandable). Cette tension est num\'eris\'ee par un convertisseur analogique-num\'erique de plage d'entr\'ee $[\SI{-5}{\volt}\,;\,\SI{+5}{\volt}]$.

\begin{enumerate}[leftmargin=1.5em]
  \item Rappeler le principe de fonctionnement d'un r\'esolveur. Expliquer, \`a l'aide d'un sch\'ema, pourquoi ce capteur d\'elivre le cosinus et le sinus de la position angulaire du moteur, et pr\'eciser ce qui le distingue d'un codeur incr\'emental.
  \item Montrer que le pas angulaire commandable vaut $\delta\theta = 1{,}72\cdot10^{-5}\,^\circ$, puis \'etablir la valeur de $\delta v$.
  \item Le capteur est consid\'er\'e comme suffisamment pr\'ecis si le pas de quantification $q$ du CAN est dix fois plus petit que $\delta v$. D\'eterminer le nombre de bits de codage n\'ecessaire au CAN.
\end{enumerate}

}{
\textbf{Corrig\'e.}

\begin{enumerate}[leftmargin=1.5em]
  \item Un r\'esolveur est un petit transformateur tournant. Un enroulement port\'e par le rotor est excit\'e ; le flux qu'il cr\'ee est capt\'e par \textbf{deux enroulements statoriques dispos\'es \`a $90^\circ$ l'un de l'autre}. Le champ magn\'etique du rotor \'etant inclin\'e de l'angle $\theta$ par rapport aux axes du stator, sa projection sur chacun des deux axes perpendiculaires fait appara\^itre naturellement un $\cos\theta$ et un $\sin\theta$ :

\begin{center}
\begin{tikzpicture}[scale=1.25, >=Stealth]
  % stator : deux axes perpendiculaires
  \draw[gray!60] (-1.5,0) -- (1.5,0);
  \draw[gray!60] (0,-1.5) -- (0,1.5);
  \node[gray!70, right] at (1.5,0) {\footnotesize bobine $\cos$};
  \node[gray!70, above] at (0,1.5) {\footnotesize bobine $\sin$};
  \draw[gray!40] (0,0) circle (1.3);
  % rotor incline
  \draw[-latex, very thick, colExo] (0,0) -- (55:1.15) node[above right]{\small $\vec{B}$ (rotor)};
  % projections
  \draw[dashed, red!70] (55:1.15) -- (0.6596,0);
  \draw[dashed, red!70] (55:1.15) -- (0,0.9420);
  \draw[-latex, thick, red!80] (0,0) -- (0.6596,0) node[below, pos=0.55]{\footnotesize $\propto\cos\theta$};
  \draw[-latex, thick, red!80] (0,0) -- (0,0.9420) node[left, pos=0.6]{\footnotesize $\propto\sin\theta$};
  \draw (0.42,0) arc (0:55:0.42);
  \node at (0.40,0.20) {\small $\theta$};
\end{tikzpicture}
\end{center}

  Le rapport des deux tensions donne $\tan\theta$, dont l'arc-tangente restitue $\theta$ sans ambigu\"it\'e sur un tour. Contrairement \`a un codeur incr\'emental, qui ne compte que des impulsions et exige une prise d'origine \`a chaque mise sous tension, le r\'esolveur est un capteur \textbf{absolu} : la position est connue d\`es la mise sous tension. Il est en outre purement \'electromagn\'etique (pas d'\'electronique ni d'optique dans le capteur), donc robuste aux vibrations, \`a la temp\'erature et aux perturbations \'electromagn\'etiques -- d'o\`u son emploi syst\'ematique sur les servomoteurs brushless industriels, o\`u l'autopilotage exige de conna\^itre la position du rotor \`a tout instant.

  \item Le mot de commande est cod\'e sur $24$ bits, soit $N = 2^{24}$ valeurs possibles r\'eparties sur le d\'ebattement de $288^\circ$ :
  \[ \delta\theta = \dfrac{288}{2^{24}} = \dfrac{288}{16\,777\,216} \approx 1{,}72\cdot10^{-5}\,^\circ \]
  La tension image \'evolue lin\'eairement de $\SI{-5}{\volt}$ \`a $\SI{+5}{\volt}$, soit une excursion de $\SI{10}{\volt}$ pour les $288^\circ$ de d\'ebattement. La variation de tension correspondant \`a un pas angulaire est donc :
  \[ \delta v = \dfrac{10}{288}\times\delta\theta = \dfrac{10}{288}\times 1{,}72\cdot10^{-5} \approx \SI{5,96e-7}{\volt} \]
  soit environ $\SI{0,6}{\micro\volt}$ -- autrement dit $\delta v = 10/2^{24}$, ce qui \'etait pr\'evisible : les $24$ bits de commande d\'ecoupent directement l'excursion de $\SI{10}{\volt}$.

  \item Le crit\`ere impose $q = \delta v/10 \approx 6\cdot10^{-8}\,\si{\volt}$. Le nombre de segments \`a distinguer sur la plage $[\SI{-5}{\volt}\,;\,\SI{+5}{\volt}]$ vaut :
  \[ N_{seg} = \dfrac{10}{q} = \dfrac{10}{6\cdot10^{-8}} \approx 1{,}6\cdot10^{8} \]
  \[ N_{bits} = \left\lceil \log_2\left(1{,}6\cdot10^{8}\right)\right\rceil = \lceil 27{,}25 \rceil = \textbf{28 bits} \]
\end{enumerate}

\textbf{Conclusion.} Il faudrait un CAN de $28$ bits pour que la cha\^ine d'acquisition du r\'esolveur ne d\'egrade pas la r\'esolution de commande de $24$ bits : c'est bien au-del\`a des CAN industriels courants ($12$ \`a $16$ bits, $24$ bits pour les convertisseurs $\Sigma\Delta$ de mesure lente). Le maillon limitant de l'asservissement de position n'est donc pas le servomoteur brushless ni le r\'esolveur, mais la \textbf{num\'erisation de la mesure} -- il faut soit r\'eduire la plage de tension \`a num\'eriser autour du point de fonctionnement, soit accepter une pr\'ecision de positionnement moindre que le pas de commande th\'eorique.
}
\end{tcolorbox}

\begin{tcolorbox}[exoBox,
  title={Exercice 4 -- \difficulte{3} \hfill\normalfont\footnotesize 2024, \'Epreuve \'ecrite disciplinaire appliqu\'ee (Urbanloop), questions 48, 49, 51, 52 et 53}
]
\CAPETSHOW{
Le bo\^itier d'alimentation qui alimente le bus continu $\SI{72}{\volt}$ des capsules Urbanloop repose sur un pont redresseur PD2 (pont de Gr\ae tz, 4 diodes suppos\'ees parfaites), aliment\'e par le r\'eseau EDF monophas\'e $v_1(t)$ et d\'elivrant une tension $u_S(t)$ c\^ot\'e continu sur une charge $R_{ch}$.

\begin{center}
\begin{circuitikz}[european, scale=0.8]
\draw
(0,2) node[circ]{} to[sV, l=$v_1$] (5,2) node[circ]{}
;
\draw (0,2) to[D, l_=$D_1$] (2.5,4.5) node[circ]{};
\draw (5,2) to[D, l_=$D_2$] (2.5,4.5);
\draw (2.5,-0.5) node[circ]{} to[D, l_=$D_3$] (0,2);
\draw (2.5,-0.5) to[D, l_=$D_4$] (5,2);
\draw (2.5,4.5) to[short] (2.5,5.3) -- (8,5.3) to[R, l_=$R_{ch}$] (8,-1.3) -- (2.5,-1.3) to[short] (2.5,-0.5);
\node[above right] at (2.6,5.15) {$+$};
\node[below right] at (2.6,-1.15) {$-$};
\draw[-latex] (9.5,5.3) -- (9.5,-1.3) node[midway,right]{$u_S(t)$};
\end{circuitikz}
\end{center}

\begin{enumerate}[leftmargin=1.5em]
  \item Indiquer si le bo\^itier d'alimentation (pont PD2) est r\'eversible en tension et en courant, du point de vue du bus continu. Justifier \`a l'aide du sch\'ema.
  \item Lors du freinage de la capsule, la machine synchrone fonctionne en g\'en\'eratrice et cherche \`a renvoyer de l'\'energie vers le bo\^itier d'alimentation. Conclure sur les risques encourus si aucune solution n'est pr\'evue pour absorber cette \'energie.
  \item Pour stocker cette \'energie de freinage, on utilise des supercondensateurs de caract\'eristiques $\SI{65}{\farad}$, $\SI{16,2}{\volt}$. D\'eterminer l'agencement \`a retenir et le nombre minimal $N_{sc}$ de supercondensateurs pour supporter la tension du bus continu de $\SI{72}{\volt}$, puis la capacit\'e \'equivalente $C_{eq}$ et la tension maximale $U_{Ceq,max}$ de l'association.
  \item L'\'energie cin\'etique r\'ecup\'erable lors d'un freinage complet depuis $V_0$ vaut $\Delta E_C = \SI{-139}{\kilo\joule}$, dont $\SI{80}{\percent}$ seulement sont effectivement r\'ecup\'er\'es. Au d\'ebut du freinage, la tension aux bornes du condensateur \'equivalent vaut $U_{Ceq}(t_c) = U_0 = \SI{72}{\volt}$. Calculer la tension $U_{Ceq,f}$ atteinte \`a la fin de la phase de freinage.
  \item Conclure sur l'agencement propos\'e \`a la question~3. Proposer, le cas \'ech\'eant, des solutions pour \'eviter la destruction des supercondensateurs.
\end{enumerate}

}{
\textbf{Corrig\'e.}

\begin{enumerate}[leftmargin=1.5em]
  \item D'apr\`es le sch\'ema, quel que soit le signe de $v_1(t)$, le courant dans $R_{ch}$ ne peut circuler que dans un seul sens (impos\'e par l'orientation des diodes $D_1$ \`a $D_4$) et la tension $u_S(t)$ ne peut \^etre que positive (ou nulle). Le pont PD2, non command\'e, ne fonctionne donc que dans un seul quadrant du plan $(u_S, i_S)$~: il est \textbf{unidirectionnel en courant et en tension}, donc non r\'eversible, ce que confirme le corrig\'e officiel (\guillemotleft\ comme le pont PD2 est unidirectionnel en courant et en tension, les bo\^itiers sont unidirectionnels en courant et en tension \guillemotright).

  \item Pendant le freinage, de l'\'energie est renvoy\'ee vers les bo\^itiers d'alimentation. Le pont PD2 ne pouvant pas absorber un courant invers\'e, cette \'energie ne peut pas transiter vers le r\'eseau~: si elle est trop importante et que rien n'est pr\'evu pour l'absorber, elle risque d'endommager les bo\^itiers d'alimentation, voire d'entra\^iner leur destruction.

  \item Pour supporter une tension sup\'erieure \`a leur tension nominale, les supercondensateurs doivent \^etre associ\'es \textbf{en s\'erie}. Pour que les $\SI{72}{\volt}$ se r\'epartissent \'equitablement sur $N_{sc}$ \'el\'ements de $\SI{16,2}{\volt}$ :
  \[ 16{,}2 \times N_{sc} \geq 72 \;\Longrightarrow\; N_{sc} \geq 4{,}4 \;\Longrightarrow\; N_{sc} = 5 \text{ (minimum)} \]
  \[ C_{eq} = \dfrac{65}{5} = \SI{13}{\farad} \qquad\text{et}\qquad U_{Ceq,max} = 5 \times 16{,}2 = \SI{81}{\volt} \]

  \item L'\'energie effectivement stock\'ee dans le condensateur \'equivalent est $\Delta E_{Cond} = -0{,}8\,\Delta E_C = 0{,}8\times 139 = \SI{111,2}{\kilo\joule}$. Or, pour un condensateur :
  \[ \Delta E_{Cond} = \tfrac{1}{2}C_{eq}\left(U_{Ceq,f}^2 - U_0^2\right) \;\Longrightarrow\; U_{Ceq,f} = \sqrt{U_0^2 + \dfrac{2\,\Delta E_{Cond}}{C_{eq}}} \]
  \[ U_{Ceq,f} = \sqrt{72^2 + \dfrac{2\times 111\,200}{13}} = \sqrt{5\,184 + 17\,108} = \sqrt{22\,292} \approx \SI{149}{\volt} \]

  \item L'agencement propos\'e \`a la question~3 \textbf{ne convient pas}~: on aurait $U_{Ceq,f} \approx \SI{149}{\volt} > U_{Ceq,max} = \SI{81}{\volt}$, chaque cellule voyant alors $149/5 \approx \SI{29,8}{\volt}$, soit pr\`es du double de sa tension nominale de $\SI{16,2}{\volt}$ -- destruction assur\'ee. Les solutions possibles sont :
  \begin{itemize}[leftmargin=1.5em]
    \item associer \textbf{dix} supercondensateurs $\SI{65}{\farad}$ / $\SI{16,2}{\volt}$ en s\'erie ($U_{Ceq,max} = \SI{162}{\volt}$) ;
    \item prendre des supercondensateurs de caract\'eristiques diff\'erentes (capacit\'e ou tension nominale plus \'elev\'ee) ;
    \item mettre en place un syst\`eme de dissipation d'\'energie \`a base de r\'esistances (\emph{rh\'eostat de freinage}) pour \'evacuer l'exc\'edent.
  \end{itemize}
\end{enumerate}

\textbf{Conclusion.} Le pi\`ege de cette partie est de s'arr\^eter \`a la question~3 : $N_{sc} = 5$ satisfait bien le crit\`ere de \emph{tension de bus} ($\SI{81}{\volt} > \SI{72}{\volt}$), mais pas le crit\`ere \emph{\'energ\'etique}. Dimensionner un stockage de freinage impose de v\'erifier les deux : la tension maximale support\'ee \emph{et} l'\'energie \`a absorber. Ici c'est l'\'energie qui commande, et elle exige de doubler le nombre d'\'el\'ements.
}
\end{tcolorbox}

\begin{tcolorbox}[exoBox,
  title={Exercice 5 -- \difficulte{3} \hfill\normalfont\footnotesize exercice in\'edit}
]
\CAPETSHOW{
Une machine synchrone triphas\'ee \`a rotor bobin\'e (mod\`ele de Behn-Eschenburg, r\'esistance statorique n\'eglig\'ee) fonctionne en moteur, aliment\'ee sous tension simple efficace $V_s = \SI{230}{\volt}$. Elle absorbe un courant de ligne $I_s = \SI{100}{\ampere}$ avec un facteur de puissance $\cos\varphi = 0{,}85$ inductif ($\varphi > 0$, courant en retard sur la tension). La r\'eactance synchrone est $X_s = \SI{2}{\ohm}$.

\begin{center}
\begin{tikzpicture}[scale=1.15, >=Stealth]

  \coordinate (O)  at (0,0);
  \coordinate (Vs) at (2.30,0);
  \coordinate (Is) at (0.893,-0.553);
  \coordinate (E)  at (1.246,-1.700);

  \draw[-latex, thick, blue]    (O)  -- (Vs) node[above right]{$\underline{V}_s$};
  \draw[-latex, thick, red]     (O)  -- (Is) node[below]{$\underline{I}_s$};
  \draw[-latex, thick, teal]    (O)  -- (E)  node[below]{$\underline{E}$};
  \draw[-latex, dashed, orange] (Vs) -- (E)  node[midway, right]{$-jX_s\underline{I}_s$};

  % parallele a I_s issue de la pointe de V_s : materialise l'angle droit du triangle de Kapp
  \draw[dotted, gray] (Vs) -- (2.938,-0.395);
  \draw[gray] (2.487,-0.116) -- (2.371,-0.303) -- (2.184,-0.187);

  \draw (0.85,0) arc (0:-31.8:0.85);
  \node at (0.980,-0.281) {\small $\varphi$};

  \draw (1.60,0) arc (0:-53.8:1.60);
  \node at (1.378,-0.343) {\small $\delta$};

\end{tikzpicture}

{\footnotesize Diagramme de Kapp en convention moteur~: m\^eme \'echelle de tension pour $\underline{V}_s$, $\underline{E}$ et $X_s\underline{I}_s$~; $-jX_s\underline{I}_s$ est perpendiculaire \`a $\underline{I}_s$.}
\end{center}

\begin{enumerate}[leftmargin=1.5em]
  \item Rappeler l'\'equation vectorielle du diagramme de Kapp reliant $\underline{V}_s$, $\underline{E}$, $X_s$ et $\underline{I}_s$ (r\'esistance statorique n\'eglig\'ee), en convention moteur.
  \item En projetant cette \'equation sur les axes horizontal et vertical (axe horizontal port\'e par $\underline{V}_s$), calculer la force \'electromotrice $E$ et l'angle de charge $\delta$.
  \item Rappeler la signification physique de l'angle $\delta$ et la condition de d\'ecrochage de la machine synchrone \`a vitesse impos\'ee par le r\'eseau.
\end{enumerate}

}{
\textbf{Corrig\'e.}

\begin{enumerate}[leftmargin=1.5em]
  \item En convention moteur, r\'esistance statorique n\'eglig\'ee :
  \[ \underline{V}_s = \underline{E} + jX_s\underline{I}_s \quad\Longleftrightarrow\quad \underline{E} = \underline{V}_s - jX_s\underline{I}_s \]
  \item Avec $\underline{V}_s$ pris comme origine des phases ($\underline{V}_s = V_s$) et $\underline{I}_s = I_s(\cos\varphi - j\sin\varphi)$ ($\varphi$ en retard) :
  \[ jX_s\underline{I}_s = jX_sI_s\cos\varphi + X_sI_s\sin\varphi \]
  \[ \underline{E} = V_s - X_sI_s\sin\varphi - jX_sI_s\cos\varphi \]
  Calcul de $\varphi$ : $\varphi = \arccos(0{,}85) = 31{,}8^\circ$, donc $\sin\varphi = 0{,}527$.
  \[ E_x = V_s - X_sI_s\sin\varphi = 230 - 2\times100\times0{,}527 = 230 - 105{,}4 = \SI{124,6}{\volt} \]
  \[ E_y = -X_sI_s\cos\varphi = -2\times100\times0{,}85 = \SI{-170}{\volt} \]
  \[ E = \sqrt{E_x^2 + E_y^2} = \sqrt{124{,}6^2 + 170^2} = \sqrt{15\,525 + 28\,900} = \sqrt{44\,425} \approx \SI{210,8}{\volt} \]
  \[ \delta = \arctan\left(\dfrac{|E_y|}{E_x}\right) = \arctan\left(\dfrac{170}{124{,}6}\right) \approx 53{,}8^\circ \]
  Le signe \textbf{n\'egatif} de $E_y$ signifie que la phase de $\underline{E}$ est n\'egative ($\underline{V}_s$ \'etant l'origine des phases)~: $\underline{E}$ -- donc le rotor -- est \textbf{en retard} de $53{,}8^\circ$ sur $\underline{V}_s$, ce qui est bien le sens attendu en fonctionnement moteur. Dans le diagramme de Kapp, $\underline{E}$ se place donc \emph{sous} l'axe de $\underline{V}_s$~; l'angle de charge est compt\'e en valeur absolue, $\delta = 53{,}8^\circ > 0$.
  \item L'angle de charge $\delta$ repr\'esente le d\'ephasage entre la f.\'e.m.\ $\underline{E}$ (li\'ee \`a la position du rotor) et la tension statorique $\underline{V}_s$ (li\'ee au champ tournant)~: en \textbf{moteur}, c'est le champ tournant statorique qui entra\^ine le rotor, lequel est donc \textbf{en retard} sur lui -- $\underline{E}$ est en retard sur $\underline{V}_s$, ce que traduit ici $E_y < 0$~; en \textbf{g\'en\'eratrice} (alternateur), c'est au contraire la machine d'entra\^inement qui pousse le rotor, qui prend alors de l'\textbf{avance} sur le champ tournant, et $\underline{E}$ passe en avance sur $\underline{V}_s$. Le signe de $\delta$ change donc avec le sens du transfert de puissance. La puissance \'electromagn\'etique s'\'ecrit $P_{em} = 3\dfrac{V_sE}{X_s}\sin\delta$ ($\delta$ pris en valeur absolue), maximale pour $\delta = 90^\circ$. Si la charge impose un couple sup\'erieur au couple maximal $C_{max} = 3V_sE/(X_s\Omega)$ atteint pour $\delta = 90^\circ$, la machine \textbf{d\'ecroche}~: elle perd le synchronisme et s'arr\^ete. Ce risque est \'elimin\'e en autopilotage (PMSM), o\`u la fr\'equence d'alimentation est asservie \`a la position du rotor.
\end{enumerate}

\textbf{Conclusion.} Au point de fonctionnement \'etudi\'e, $\delta = 53{,}8^\circ$ repr\'esente d\'ej\`a $\SI{60}{\percent}$ de l'angle de d\'ecrochage ($90^\circ$)~: la r\'eserve de couple avant perte de synchronisme n'est que de $1/\sin(53{,}8^\circ) \approx 1{,}24$, soit $\SI{24}{\percent}$. C'est une marge faible pour une machine \`a fr\'equence impos\'ee par le r\'eseau~; c'est pr\'ecis\'ement ce risque que supprime l'autopilotage des PMSM de traction, o\`u $\delta$ est maintenu \`a sa valeur de consigne quelle que soit la charge.
}
\end{tcolorbox}

\CAPETENDFILE
